\subsection{Cross-Representation Consistency}
\label{sec:cross_representation}

A scientific simulation admits several representations of the same underlying
physical system. A PDE can be specified mathematically, described in natural
language, implemented as source code, and observed through a numerical trajectory.
These four representations are not paraphrases of one another: each is produced by a
different process, each can be checked independently, and each constitutes a
ground-truth description of the same physical system. This provides a controlled
setting for asking whether LLMs can integrate evidence across representations, and,
when the representations conflict, which one they treat as authoritative.

\paragraph{Task.}
We build items from $32$ solvers, eight from each of four PDE families (heat, wave,
Burgers, Navier--Stokes); $16$ are human-written and $16$ synthetic. For each solver
we construct four views of the same physical system. In the \emph{consistent}
condition, all four views describe the same system. In each \emph{inconsistent}
condition, exactly one view is corrupted while the remaining three agree.

\emph{Views:} (1) source code, with all comments and docstrings removed; (2) a
numerical trajectory, rendered as $10$ time samples spanning the full run and
resampled onto a common uniform normalized spatial grid ($\leq 64$ probes for 1-D
fields, $\approx 14 \times 14$ for 2-D), printed in a single fixed \texttt{\%+.4e}
format; (3) a natural-language description of the physical process (median $92$
words); and (4) a mathematical specification of the governing PDE in \LaTeX,
including parameters, boundary conditions, and initial conditions.

\emph{Corruptions:} We corrupt code with a human-authored edit that breaks the
physics of the discretization while leaving the program runnable --- a sign flip in
a flux or update term, a reversed difference direction, or an altered coefficient
(median $1$ changed non-blank line of a median $40$-line solver, range $1$--$9$).
Every code corruption was verified by execution and carries a recorded failure mode
(blow-up to NaN, spurious oscillations, negative temperatures). We corrupt the
natural-language description by altering the claimed physical behaviour --- an
absorbing boundary in place of a reflecting one, a smooth decay in place of shock
formation, a clamped boundary in place of a flux boundary --- which propagates
through the remainder of the narrative (median $47$ of $92$ words changed, first
divergence a median $37\%$ into the text). We corrupt the mathematical specification
with a minimal symbolic edit, a median of a single token: the sign of a term or of a
declared parameter ($\partial_{tt}u = -\nabla\!\cdot\!(c^2\nabla u)$; $\beta \to
-\beta$; a reversed flux divergence).

\emph{Trajectory ladder:} We introduce four types of trajectory corruption, ordered
from structurally empty to physically subtle: an execution-based invalid trajectory
(the corrupted solver's own output, decimated to the same frames); a trajectory
drawn from another system (coherent physics, wrong referent); random values matched
in shape, mean, and variance; and values shuffled within the original trajectory,
which preserves the shape and the entire value histogram exactly, so that nothing
but the spatiotemporal arrangement distinguishes it. Because all four rungs corrupt
the \emph{same} view and are rendered through the same grid and format, they are
mutually comparable in a way the across-view conditions are not.

\emph{Lexical ablation:} We additionally vary whether source-code variables retain
their original names or are replaced with obfuscated identifiers (an AST rewrite
mapping every user-defined name to \texttt{foobar\_$k$}, preserving semantics and
library calls), allowing us to test whether judgments depend on lexical cues in the
code.

\emph{Response format:} We present the four views in randomized order behind a
neutral legend, with the corrupted view counterbalanced so that it occupies each
slot equally often, and ask the model to produce a structured response containing:
(i) whether the views agree, (ii) the identity of the inconsistent view if they do
not, (iii) the physical system represented by the majority of the views (PDE class
and numerical method), and (iv) a free-text justification. Prompts are listed in
Table~\ref{tab:xr_prompt}; the full template is given in
Appendix~\ref{app:xr_prompt}.

\begin{table}[h]
\centering
\small
\begin{tabular}{@{}p{0.14\linewidth}p{0.78\linewidth}@{}}
\toprule
Component & Text \\
\midrule
Instruction & Below are four representations of a physical system, labelled View 1
to View 4. \dots\ Do these representations all describe the same physical system?
Either all four representations are in agreement, or exactly one of them is
inconsistent with the other three. \\
Legend & View $i$: \{Python source code for a numerical solver $\mid$ the governing
equation, in LaTeX $\mid$ the numerical solution field produced by running a solver
$\mid$ a natural-language description of the physical process\} \\
Response & Answer with a JSON object containing exactly these fields:
\texttt{agree}, \texttt{outlier}, \texttt{system\_pde\_class},
\texttt{system\_num\_method}, \texttt{justification}. Output only the JSON object. \\
\bottomrule
\end{tabular}
\caption{Prompt components. The legend assigns one line per slot in that item's
permuted order, and the four view bodies follow, wrapped in \texttt{<view\_$i$>}
tags. The prompt never names a PDE class or numerical method, since the model is
asked to recover both.}
\label{tab:xr_prompt}
\end{table}

\paragraph{Experimental design.}
We conduct experiments over eight open-weight reasoning models spanning
2025-01 to 2026-08 and four model families, held to a narrow $27.8$--$32.8$B
parameter band so that release date is not confounded with scale. All models used
are listed in Appendix~\ref{app:xr_models}. Every model is run reasoning-enabled,
with reasoning requested through the chat template rather than the prompt text, so
that reasoning is held constant and the comparison across models is clean. Crossing
$32$ solvers $\times$ $8$ conditions $\times$ $2$ naming levels $\times$ $2$ order
permutations gives $1{,}024$ items per model. We draw $k = 3$ samples per item under
one uniform decoding regime applied to every model; decoding settings are given in
Appendix~\ref{app:xr_inference}.

\emph{Metric:} Because the design is deliberately $7{:}1$ corrupted-to-clean, raw
accuracy is not interpretable --- a model that always answers ``disagree'' scores
$87.5\%$. We therefore report the flag rate on corrupted items against the
false-alarm rate on consistent items, and summarise the pair with the bias-corrected
detection sensitivity $d'$.

\emph{Uncertainty:} We report $95\%$ confidence intervals bootstrapped over the $32$
solvers rather than over the $1{,}024$ items, since items sharing a solver are not
independent; the effective sample size is $32$, not $1{,}024$. Naming contrasts are
computed within-solver, matched on (solver, condition, reasoning, model) before
differencing.

\subsubsection*{Results}

\paragraph{Detection depends far more on how a system is corrupted than on whether
it is corrupted.}
Against a false-alarm rate of $0.377$ $[0.297, 0.457]$ on fully consistent items,
detection varies by more than a factor of two across corruptions. Within the
trajectory ladder --- the one comparison that is severity-controlled, since all four
rungs corrupt the same view --- the ordering is the opposite of what a numerically
attentive reader would produce. The invalid solver's own output is caught most often
($0.844$ $[0.773, 0.906]$), followed by another system's trajectory ($0.635$
$[0.559, 0.711]$), while a field of pure shape-matched noise is caught only $0.555$
$[0.465, 0.645]$ of the time and a within-trajectory permutation --- which preserves
every marginal statistic and differs from the valid field only in its arrangement
--- only $0.521$ $[0.445, 0.600]$, barely above the false-alarm floor. Restricting
to the $15$ of $32$ invalid solvers whose output remains finite drops the
execution-based condition from $0.952$ to $0.721$, still the highest of any
condition but a considerably smaller margin. Across views, the corrupted equation is
flagged most often ($0.686$ $[0.594, 0.773]$), the corrupted description next
($0.602$ $[0.520, 0.686]$), and the corrupted code least ($0.469$ $[0.389, 0.551]$);
we report this ordering but do not read it as a ranking of trust, because the three
corruptions are not severity-matched. \textbf{Takeaway:} models do not appear to
read the numerical field as a physical field --- a structureless array survives
scrutiny nearly as often as a consistent one, and the invalid solver's output is
caught largely because it is visibly pathological.

\paragraph{Errors reveal a systematic blame bias against executable representations.}
When a model does flag an item, it names the correct view $67.0\%$ $[61.8, 72.4]$ of
the time --- above the best fixed strategy, so blame does track the real outlier ---
but the residual errors are far from uniform. Defining the false-blame rate for a
view as $P(\text{accused} = m \mid \text{true outlier} \neq m)$, code is accused
$15.8\%$ $[11.7, 20.1]$ of the time when it is not the broken view and the
trajectory $11.9\%$ $[9.3, 14.6]$, against $2.3\%$ $[0.9, 4.1]$ for the description
and $1.1\%$ $[0.6, 1.6]$ for the equation --- a fourteen-fold gap between the most-
and least-suspected representation. The pattern is sharpest where there is nothing
to find: on the consistent items that were falsely flagged, $48.7\%$ of accusations
fall on the code and $39.9\%$ on the trajectory, leaving $11.4\%$ for the
description and the equation combined. Among all misattributions on corrupted items,
$64.5\%$ name the code. This bias runs precisely opposite to accuracy: the code is
the view least often identified when it genuinely is the outlier ($30.9\%$ of
code-corrupted items, versus $49.4\%$ for the equation), and simultaneously the view
most often accused when it is not. \textbf{Takeaway:} models treat the equation and
the prose as the reference frame and adjudicate the executable artifacts against
them --- the inverse of the ordering an auditor of scientific software would want,
since the code and its output are what determine the computation actually performed.

\paragraph{Removing lexical cues changes confidence, not competence.}
Obfuscating identifiers does not measurably impair the ability to catch broken code.
In a within-solver paired comparison, the end-to-end rate at which a model both
flags and correctly names genuinely broken code changes by $-2.3$pp $[-9.8, +5.5]$,
and localization given a flag by $+1.6$pp $[-8.3, +12.7]$; neither interval excludes
zero. Nor does stripped code draw more suspicion when it is innocent: false blame on
code moves $+0.8$pp $[-2.1, +3.9]$. Identification of the numerical method, which is
carried only by the code view, is likewise flat ($0.467$ named vs.\ $0.485$
obfuscated). What obfuscation does change is the willingness to commit: models
answer ``all four agree'' $+4.8$pp $[+2.0, +7.7]$ more often and flag corrupted
items $-4.9$pp $[-7.9, -1.8]$ less often, and the blame that disappears comes
disproportionately off the trajectory ($-4.7$pp $[-8.1, -0.8]$).
\textbf{Takeaway:} removing variable names removes confidence rather than evidence,
which is inconsistent with the hypothesis that apparent cross-modal reasoning is
lexical matching between the code and the prose.

\paragraph{Reasoning shifts the decision criterion rather than the underlying
sensitivity.}
The one model run both ways isolates the effect of test-time reasoning. Enabling it
raises the flag rate on corrupted items from $0.408$ to $0.636$, but raises the
false-alarm rate on consistent items from $0.125$ to $0.375$ at the same time; $d'$
changes by $-0.241$ $[-0.602, +0.034]$. Across all four arms, $d'$ spans only
$0.52$--$0.91$ while the false-alarm rate spans $0.13$--$0.70$. Reported as accuracy
on this $7{:}1$ design, the most prolific flagger (\texttt{QwQ-32B}, $0.915$ flag
rate on corrupted items) would appear to be the strongest model despite a $0.703$
false-alarm rate. \textbf{Takeaway:} what separates these models is almost entirely
how readily they declare a disagreement, not how well they tell disagreement apart
from agreement, and a bias-corrected measure is necessary to see this.

\paragraph{Summary.}
Given four views of one physical system and told that at most one is wrong, current
$32$B reasoning models detect the inconsistency at rates between $0.47$ and $0.84$
depending on which view was corrupted, against a $0.38$ base rate of flagging
systems that are entirely consistent. Their failures are structured rather than
noisy: they under-read numerical fields, they treat the equation and the prose as
authoritative and the code and trajectory as suspect, and the discrimination they do
show is insensitive to variable naming but highly sensitive to decoding-time
reasoning, which moves the decision criterion without improving the underlying
signal.

% ===========================================================================
%  APPENDIX
% ===========================================================================
\section{Cross-Representation Consistency}
\label{app:cross_representation}

\subsection{Models and Prompts}
\label{app:xr_models_prompts}

\subsubsection{Models}
\label{app:xr_models}


Table~\ref{tab:xr_models} presents the corresponding model versions for the LLMs
used in our cross-representation experiments. All models are open-weight and were
imported from HuggingFace and served locally.

\begin{table}[h]
\centering
\small
\begin{tabular}{@{}ll@{}}
\toprule
LLM & Model Version  \\
\midrule
R1-Distill-32B      & \texttt{deepseek-ai/DeepSeek-R1-Distill-Qwen-32B}    \\
QwQ-32B             & \texttt{Qwen/QwQ-32B}                                \\
Qwen3-32B           & \texttt{Qwen/Qwen3-32B}                               \\
Nemotron-3-Nano-30B & \texttt{nvidia/NVIDIA-Nemotron-3-Nano-30B-A3B-BF16}  \\
GLM-4.7-Flash       & \texttt{zai-org/GLM-4.7-Flash}                        \\
Qwen3.5-27B         & \texttt{Qwen/Qwen3.5-27B}                            \\
Qwen3.6-27B         & \texttt{Qwen/Qwen3.6-27B}                             \\
Qwen3.8-27B         & \texttt{Qwen/Qwen3.8-27B}                             \\
\bottomrule
\end{tabular}
\caption{LLMs and model versions used in the experiments. Every model is run reasoning-enabled.}
\label{tab:xr_models}
\end{table}

\subsubsection{Inference configuration}
\label{app:xr_inference}

Measured with each model's own tokenizer and chat template, prompts run to a median
of $7{,}472$--$9{,}073$ tokens depending on the model, and a maximum of $33{,}097$;
the long tail comes from 2-D fields, whose rendered trajectory tables dominate the
prompt.

We use the following sampling parameters for all models: $T = 0.6$, $p = 0.95$,
top-$k = 20$. For each prompt, we sample $k = 3$ generations per model, each retained
as a separate observation. The context window is set to each model's full declared
context, capped at $69{,}632$ tokens. We use sampling rather than greedy decoding
because every checkpoint in the roster ships \texttt{do\_sample=True} and the
reasoning-model cards state that greedy decoding induces repetition.

\subsubsection{Prompt}
\label{app:xr_prompt}

The following prompt was used for the consistency experiment.

\begin{verbatim}
Below are four representations of a physical system, labelled View 1 to View 4.

Legend:
{legend}

{views}

Do these representations all describe the same physical system? Either all four
representations are in agreement, or exactly one of them is inconsistent with the
other three.

Answer with a JSON object containing exactly these fields:
- "agree": "yes" if all four agree, "no" if one is inconsistent.
- "outlier": which view is inconsistent, as "view_1", "view_2", "view_3" or
  "view_4". Use "none" if and only if "agree" is "yes".
- "system_pde_class": the class of partial differential equation that the majority
  of the views describe.
- "system_num_method": the numerical method that the majority of the views use.
- "justification": what specifically is inconsistent, or why you judge the views
  consistent.

Output only the JSON object.
\end{verbatim}

\texttt{\{legend\}} is replaced by one line per slot, \texttt{View $i$:} followed by
the corresponding string below; \texttt{\{views\}} by the four view bodies, each
wrapped in \texttt{<view\_$i$>} tags. The legend strings are fixed, one per
representation, and state what each view \emph{is} in a uniform register without
indicating which is more reliable:

\begin{verbatim}
Python source code for a numerical solver
the governing equation, in LaTeX
the numerical solution field produced by running a solver
a natural-language description of the physical process
\end{verbatim}

Slot order is permuted
per item, with the corrupted view counterbalanced across slots, so that position
carries no information about which view is the outlier.

Trajectories are rendered with a fixed header stating what the axes mean and
explicitly disclaiming units. The stored trajectory is a bare array with no
accompanying axis metadata --- no coordinates, no time step, no grid spacing --- so
there is no basis for labelling its axes in physical units, and the header states
normalized coordinates instead. Those quantities are not absent from the item as a
whole: the code defines a time step in $26$ of $32$ solvers and a grid spacing in
$25$, and all $32$ mathematical specifications state a spatial domain. They are
simply not attached to the numerical field, which is the view being rendered.

\begin{verbatim}
Numerical solution field, sampled on a uniform normalized grid.
Rows are successive time samples from the start of the run (t=0.00) to the end
(t=1.00) in equal steps.
Within a row, values run over the spatial domain from 0.00 to 1.00 in equal steps
along each axis.
No physical units, grid spacing, or time step are available for this field.

t=0.00:  +3.2500e+02  +3.2500e+02  +3.2500e+02  ...
\end{verbatim}